% ---------------------------------------------------------------------
% Proposed Method (~4.5 pages incl. TikZ figures + algorithm).
% ---------------------------------------------------------------------

\newcommand{\R}{\mathbb{R}}
\newcommand{\bone}{\mathbf{1}}

\subsection{Problem Formulation}
\label{ssec:formulation}

We frame audio deepfake detection as an \emph{open-set} binary classification problem.
Given a waveform $\mathbf{x} \in \R^{T}$ sampled at 16~kHz, a model produces a score
$s(\mathbf{x}) \in \R$ such that $s(\mathbf{x}) > \tau$ indicates AI-generated speech and
$s(\mathbf{x}) \leq \tau$ indicates human speech, where $\tau$ is a decision threshold.
Training uses a labeled set $\mathcal{D} = \{(\mathbf{x}_i, y_i)\}_{i=1}^N$ with
$y_i \in \{0, 1\}$ denoting bona-fide (0) and spoof (1). Crucially, the set of spoofing
algorithms observed at test time, $\mathcal{A}_{\text{test}}$, is not contained in the
training set $\mathcal{A}_{\text{train}}$: new generators, new languages, and new
transmission channels all shift the test distribution. The design goal is therefore not
to separate $\mathcal{A}_{\text{train}}$ from bona-fide speech maximally, but to learn a
representation in which \emph{bona-fide speech occupies a compact region} and any
synthetic artifact---seen or unseen---falls outside it.

AuralGuard implements this goal with five learned components: two parallel front-end
views (Sections~\ref{ssec:viewa} and~\ref{ssec:viewb}), a gated cross-attention fusion
module (Section~\ref{ssec:fusion}), a spectro-temporal graph-attention back-end
(Section~\ref{ssec:backend}), and a multi-part one-class objective
(Section~\ref{ssec:objective}). Figure~\ref{fig:architecture} provides an overview;
Algorithm~\ref{alg:training} summarizes training.

\begin{figure*}[t]
\centering
\resizebox{\textwidth}{!}{%
\begin{tikzpicture}[node distance=4.5mm]

% ---------------- Input ----------------
\node[io] (wav) {Waveform $\mathbf{x}\in\R^{T}$\\16~kHz};

% ---------------- View A (top branch) ----------------
\node[viewa,right=9mm of wav,yshift=9mm] (wavlm)
      {WavLM-Large\\(frozen)\\$\mathbf{H}\in\R^{L\times T'\times D}$};
\node[viewa,right=of wavlm] (latt)
      {Layer attention\\$\boldsymbol{\alpha}\!\in\!\R^{L}$, softmax\\$\bar{\mathbf{H}}=\sum_\ell\pi_\ell\mathbf{h}_\ell$};
\node[viewa,right=of latt] (projA)
      {Conv1d\\projection\\$\mathbf{F}_A\!\in\!\R^{T'\times C}$};

% ---------------- View B (bottom branch) ----------------
\node[viewb,right=9mm of wav,yshift=-9mm] (feat)
      {LFCC\,$\oplus$\,MGD\\$\oplus$\,CQT-phase\\$\mathbf{S}\!\in\!\R^{C_{\text{in}}\times F\times T''}$};
\node[viewb,right=of feat] (cnn)
      {2-D residual\\CNN ($4\times$)};
\node[viewb,right=of cnn] (projB)
      {frame feats\\$\mathbf{F}_B\!\in\!\R^{T''\times C}$};

% ---------------- Fusion ----------------
\node[fuse,right=11mm of projA,yshift=-9mm] (fusion)
      {Gated\\cross-attention\\fusion\\$\mathbf{F}\!\in\!\R^{T^*\times C}$};

% ---------------- Back-end ----------------
\node[back,right=of fusion] (aasist)
      {Spectro-temporal\\graph attention\\(AASIST, GATv2)};
\node[back,right=of aasist] (emb)
      {Utterance\\embedding\\$\mathbf{z}\in\R^{d}$};

% ---------------- Heads ----------------
\node[head,right=9mm of emb,yshift=7mm] (ocs)
      {OC-Softmax\\head $\to$ score $s$};
\node[head,right=9mm of emb,yshift=-7mm] (supcon)
      {SupCon proj.\\$g(\mathbf{z})$ (train only)};

% ---------------- Loss ----------------
\node[loss,below=7mm of supcon] (loss)
      {$\mathcal{L}=\mathcal{L}_{\text{OCS}}+\lambda_c\mathcal{L}_{\text{SupCon}}+\lambda_r\mathcal{L}_{\text{ent}}$};

% ---------------- Group boxes ----------------
\begin{scope}[on background layer]
  \node[grp,fill=blue!4,draw=blue!35,fit=(wavlm)(latt)(projA)] (grpA) {};
  \node[grp,fill=orange!4,draw=orange!35,fit=(feat)(cnn)(projB)] (grpB) {};
\end{scope}
\node[font=\scriptsize\bfseries,blue!55,above=0.5mm of grpA] {View A: Self-supervised SSL front-end};
\node[font=\scriptsize\bfseries,orange!70!black,below=0.5mm of grpB] {View B: Vocoder-artifact branch};

% ---------------- Edges ----------------
\draw[flow] (wav.east) -- ++(3mm,0) |- (wavlm.west);
\draw[flow] (wav.east) -- ++(3mm,0) |- (feat.west);
\draw[flow] (wavlm) -- (latt);
\draw[flow] (latt) -- (projA);
\draw[flow] (feat) -- (cnn);
\draw[flow] (cnn) -- (projB);
\draw[flow] (projA.east) -- ++(3mm,0) |- ([yshift=2mm]fusion.west);
\draw[flow] (projB.east) -- ++(3mm,0) |- ([yshift=-2mm]fusion.west);
\draw[flow] (fusion) -- (aasist);
\draw[flow] (aasist) -- (emb);
\draw[flow] ([yshift=2mm]emb.east) -- ++(2.5mm,0) |- (ocs.west);
\draw[flow] ([yshift=-2mm]emb.east) -- ++(2.5mm,0) |- (supcon.west);
\draw[tflow] (ocs.east) -- ([xshift=3mm]loss.east |- ocs.east) |- (loss.east);
\draw[tflow] (supcon.south) -- (supcon.south|-loss.north);

\end{tikzpicture}%
}
\caption{Overview of the AuralGuard architecture. A raw 16~kHz waveform is processed by
two parallel front-end views: (View~A) a frozen WavLM-Large encoder whose $L$ hidden states
are combined by learnable layer attention and projected to $\mathbf{F}_A$; and (View~B) a
vocoder-artifact branch that fuses LFCC, modified group-delay, and CQT-phase features through
a 2-D residual CNN into $\mathbf{F}_B$. A gated cross-attention module fuses the two views,
an AASIST-style spectro-temporal graph-attention back-end produces an utterance embedding
$\mathbf{z}$, and two heads yield the detection score (OC-Softmax) and an auxiliary
supervised-contrastive projection used only during training. Dashed red edges denote
training-time loss signals; the entropy regularizer $\mathcal{L}_{\text{ent}}$ acts on
the layer-attention weights $\boldsymbol{\pi}$ (edge omitted for clarity).}
\label{fig:architecture}
\end{figure*}

\subsection{View A: SSL Front-End with Learnable Layer Attention}
\label{ssec:viewa}

The first view extracts high-level representations from a frozen pre-trained
WavLM-Large model~\cite{chen2022wavlm}, chosen over wav2vec~2.0 and HuBERT because its
pre-training includes denoising and overlapped-speech objectives that transfer well to
artifact-sensitive tasks (we compare backbones in Section~\ref{ssec:ablation}). Given
input waveform $\mathbf{x}$, WavLM produces hidden states from $L$ transformer layers:
\begin{equation}
\mathbf{H} = [\mathbf{h}_1, \mathbf{h}_2, \ldots, \mathbf{h}_L] \in \R^{L \times T' \times D},
\end{equation}
where $L = 25$ (including the convolutional feature extractor output), $D = 1024$, and
$T'$ is the downsampled time dimension (one frame per 20~ms).

The information relevant to deepfake detection is not uniformly distributed across this
hierarchy: lower layers preserve fine acoustic and phase-adjacent structure, while upper
layers encode phonetic and semantic content~\cite{pan2024attentive}. Rather than using
only the final layer as is common~\cite{tak2022automatic}, we introduce learnable layer
attention weights $\boldsymbol{\alpha} \in \R^L$, normalized via softmax:
\begin{equation}
\bar{\mathbf{H}} = \sum_{\ell=1}^L \pi_\ell \, \mathbf{h}_\ell \in \R^{T' \times D},
\qquad
\pi_\ell = \frac{\exp(\alpha_\ell)}{\sum_{j=1}^L \exp(\alpha_j)}.
\label{eq:layerattn}
\end{equation}
An entropy regularization term
\begin{equation}
\mathcal{L}_{\text{ent}} = \sum_{\ell=1}^{L} \pi_\ell \log \pi_\ell
\label{eq:entreg}
\end{equation}
is \emph{added} to the loss (i.e., we penalize low entropy), preventing the attention
distribution from collapsing onto a single layer early in training, which we found
empirically to harm generalization: a collapsed profile overfits to the layer most
discriminative for \emph{seen} attacks. The weighted representation is projected to a
common channel dimension $C = 128$ via a 1-D convolution:
$\mathbf{F}_A = \operatorname{Conv1d}(\bar{\mathbf{H}}) \in \R^{T' \times C}$.

Freezing the backbone has three benefits: (i)~it preserves the general-purpose
representation, avoiding catastrophic overwriting by the comparatively tiny
anti-spoofing corpus; (ii)~it reduces trainable parameters from 316M to under 5M
(Table~\ref{tab:budget}); and (iii)~it stabilizes the one-class objective, which is
sensitive to representation drift. Recent evidence supports frozen SSL front-ends for
calibrated generalization~\cite{pascu2024calibrated}.

\subsection{View B: Vocoder-Artifact Branch}
\label{ssec:viewb}

SSL models are pre-trained on \emph{natural} speech for masked prediction; their
representations are therefore incentivized to discard exactly the low-level synthesis
residue---phase discontinuities, sub-band aliasing, over-smoothed spectral
envelopes---that betrays a vocoder. View~B restores this evidence explicitly
(Figure~\ref{fig:artifact}). We extract three complementary time--frequency maps:

\begin{figure}[t]
\centering
\resizebox{\columnwidth}{!}{%
\begin{tikzpicture}[node distance=3.5mm]
\node[io] (w) {$\mathbf{x}$};
\node[viewb,right=5mm of w,yshift=11mm] (lfcc) {LFCC\\60-dim, $\Delta,\Delta\Delta$};
\node[viewb,right=5mm of w] (mgd) {Modified\\group delay\\($\rho{=}0.4,\gamma{=}0.9$)};
\node[viewb,right=5mm of w,yshift=-11mm] (cqt) {CQT phase\\96 bins/oct};
\node[op,right=4mm of mgd] (cat) {$\oplus$};
\node[viewb,right=4mm of cat] (stem) {Conv $3{\times}3$, 32\\+ BN + SiLU};
\node[viewb,right=of stem] (res) {$4\times$ residual\\block (32--128)\\stride-2 freq.};
\node[viewb,right=of res] (pool) {freq.\ pooling\\$\to \mathbf{F}_B\in\R^{T''\times C}$};
\draw[flow] (w) |- (lfcc);
\draw[flow] (w) -- (mgd);
\draw[flow] (w) |- (cqt);
\draw[flow] (lfcc.east) -| (cat);
\draw[flow] (mgd) -- (cat);
\draw[flow] (cqt.east) -| (cat);
\draw[flow] (cat) -- (stem);
\draw[flow] (stem) -- (res);
\draw[flow] (res) -- (pool);
\end{tikzpicture}%
}
\caption{View B, the vocoder-artifact branch. Three complementary time--frequency
representations---magnitude cepstra (LFCC), phase-derivative structure (modified group
delay), and constant-Q phase---are stacked as input channels and encoded by a light 2-D
residual CNN into frame-level features $\mathbf{F}_B$.}
\label{fig:artifact}
\end{figure}

\begin{enumerate}
\item \textbf{LFCC:} 60-dimensional linear-frequency cepstral coefficients (512-point
  FFT, 20~linear filters, 20~ms window, 10~ms hop) with $\Delta$ and $\Delta\Delta$,
  capturing magnitude-spectrum envelope artifacts~\cite{sahidullah2015comparison}.
\item \textbf{Modified group delay (MGD):} the group-delay function
  $\tau_g(\omega) = -\,d\angle X(\omega)/d\omega$ computed with cepstral smoothing and
  compression parameters $(\rho, \gamma) = (0.4, 0.9)$, exposing phase discontinuities
  that magnitude features cannot represent.
\item \textbf{CQT phase:} octave-resolved phase from a constant-Q transform (96 bins per
  octave, 6 octaves), sensitive to the harmonic phase coherence that neural vocoders
  and codec decoders reconstruct imperfectly~\cite{todisco2017constantq,xie2025codecfake}.
\end{enumerate}

The three maps are resampled to a common $F \times T''$ grid and stacked as input
channels, $\mathbf{S} \in \R^{C_{\text{in}} \times F \times T''}$. A lightweight 2-D CNN
(a $3{\times}3$ stem followed by four residual blocks with channel widths 32--128 and
stride-2 frequency downsampling) encodes $\mathbf{S}$, and frequency-axis pooling yields
frame-level features $\mathbf{F}_B = \operatorname{CNN}_{\text{art}}(\mathbf{S}) \in
\R^{T'' \times C}$. The branch adds only 1.1M parameters. The two temporal resolutions
$T'$ and $T''$ are aligned by linear interpolation before fusion.

\subsection{Gated Cross-Attention Fusion}
\label{ssec:fusion}

Neither view dominates universally: SSL features excel on high-level inconsistencies of
codec-language-model TTS, while artifact features excel on GAN-vocoder residue
(Section~\ref{ssec:perattack}). We therefore fuse the views with a mechanism that
(i)~lets each view \emph{query} the other for corroborating evidence and (ii)~learns,
per utterance and per frame, how much to trust each view (Figure~\ref{fig:fusion}).

\begin{figure}[t]
\centering
\resizebox{0.98\columnwidth}{!}{%
\begin{tikzpicture}[x=1mm,y=1mm]
% ---- nodes (explicit coordinates; orthogonal routing) ----
\node[viewa] (fa) at (0,20) {$\mathbf{F}_A$};
\node[viewb] (fb) at (0,0)  {$\mathbf{F}_B$};
\node[fuse]  (xab) at (30,20) {MHA\\$Q{=}\mathbf{F}_A$\\$K,V{=}\mathbf{F}_B$};
\node[fuse]  (xba) at (30,0)  {MHA\\$Q{=}\mathbf{F}_B$\\$K,V{=}\mathbf{F}_A$};
\node[fuse,anchor=west] (gate) at (-7,-16)
      {gate $g_t=\sigma(\mathbf{w}_g^{\!\top}[\mathbf{F}_A;\mathbf{F}_B]_t)$};
\node[op]   (mix) at (54,10) {$+$};
\node[fuse,anchor=west] (out) at (60,10) {LN + FFN\\$\mathbf{F}\in\R^{T^*\times C}$};
% ---- straight view -> own-query MHA ----
\draw[flow] (fa) -- (xab);
\draw[flow] (fb) -- (xba);
% ---- K,V cross-links (single deliberate X) ----
\draw[flow] ([yshift=-2.2mm]fa.east) -- ([yshift=2.2mm]xba.west);
\draw[flow] ([yshift=2.2mm]fb.east) -- ([yshift=-2.2mm]xab.west);
% ---- gate inputs (outside orthogonal routes, no crossings) ----
\draw[flow] (fa.west) -- ++(-7,0) |- (gate.west);
\draw[flow] (fb.south) -- (fb.south|-gate.north);
% ---- weighted mixing ----
\draw[flow] (xab.east) -| node[pos=0.78,right,font=\tiny]{$\times\, g_t$} (mix.north);
\draw[flow] (xba.east) -- ++(4,0) |- node[pos=0.28,right,font=\tiny]{$\times (1{-}g_t)$} (mix.west);
\draw[flow] (gate.east) -| (mix.south);
\draw[flow] (mix) -- (out);
\end{tikzpicture}%
}
\caption{Gated cross-attention fusion. Each view attends to the other via multi-head
attention (MHA); a per-frame scalar gate $g_t$ mixes the two attended streams before
layer normalization and a position-wise feed-forward layer.}
\label{fig:fusion}
\end{figure}

With queries, keys, and values $\mathbf{Q}_A = \mathbf{F}_A \mathbf{W}_Q$,
$\mathbf{K}_B = \mathbf{F}_B \mathbf{W}_K$, $\mathbf{V}_B = \mathbf{F}_B \mathbf{W}_V$,
the $A{\to}B$ cross-attended representation is
\begin{equation}
\mathbf{A}_{A \rightarrow B} = \operatorname{softmax}\!\left(
    \frac{\mathbf{Q}_A \mathbf{K}_B^\top}{\sqrt{C}}\right) \mathbf{V}_B ,
\end{equation}
with the analogous $\mathbf{A}_{B \rightarrow A}$ in the reverse direction (4 heads
each). A per-frame gate
$g_t = \sigma(\mathbf{w}_g^\top [\mathbf{F}_A; \mathbf{F}_B]_t)$ produces the fused
output
\begin{equation}
\mathbf{F}_t = g_t \, \mathbf{A}_{A \rightarrow B, t} + (1 - g_t)\,
\mathbf{A}_{B \rightarrow A, t}, \qquad \mathbf{F} \in \R^{T^* \times C},
\end{equation}
followed by layer normalization and a position-wise feed-forward layer. In ablation
(Section~\ref{ssec:ablation}) this outperforms concatenation and simple averaging, and
the learned gate statistics provide a diagnostic: utterances from neural-vocoder
attacks drive $g_t$ toward View~B, whereas codec-LM attacks drive it toward View~A.

\subsection{Spectro-Temporal Graph-Attention Back-End}
\label{ssec:backend}

The fused representation $\mathbf{F}$ is processed by an AASIST-style
back-end~\cite{jung2022aasist}. Spectral nodes aggregate frequency-band information via
1-D convolutions over the channel axis and temporal nodes capture frame-level
transitions; two layers of graph attention---using the more expressive GATv2
formulation~\cite{brody2022gatv2}---exchange information within and across the two node
sets through a heterogeneous stack node. Max-graph operations and a readout layer
produce a fixed-length utterance embedding $\mathbf{z} \in \R^{d}$ with $d = 160$. We
retain this back-end over recent state-space alternatives~\cite{xiao2025xlsrmamba}
because graph attention operates naturally on our fused two-view features, whereas
published Mamba detectors are coupled to a single SSL stream; Section~\ref{ssec:sota}
nonetheless compares against them as published systems.

\subsection{One-Class Contrastive Training Objective}
\label{ssec:objective}

The complete loss combines three terms:
\begin{equation}
\mathcal{L} = \mathcal{L}_{\text{OCS}} \;+\; \lambda_c \, \mathcal{L}_{\text{SupCon}}
\;+\; \lambda_r \, \mathcal{L}_{\text{ent}} .
\label{eq:loss}
\end{equation}

\begin{figure}[t]
\centering
\begin{tikzpicture}[scale=0.85]
  % bona-fide compact region
  \draw[green!55!black,fill=green!10] (0,0) circle (0.9);
  \node[font=\scriptsize,green!40!black] at (0,-0.35) {bona fide};
  \foreach \p in {(0.2,0.3),(-0.3,0.15),(0.05,-0.05),(-0.1,0.45),(0.35,-0.2)}
    \fill[green!55!black] \p circle (1.6pt);
  % margin ring
  \draw[green!55!black,dashed] (0,0) circle (1.5);
  \node[font=\tiny,green!40!black] at (1.12,1.28) {margin $m_0$};
  % seen spoofs
  \foreach \p in {(2.6,1.1),(3.0,0.4),(2.8,-0.6),(2.3,-1.2)}
    \fill[red!65!black] \p circle (1.6pt);
  \node[font=\scriptsize,red!60!black] at (2.9,1.55) {seen spoofs};
  % unseen spoofs
  \foreach \p in {(-2.4,1.5),(-2.9,0.6),(-2.2,-1.4)}
    \draw[violet!70!black,fill=violet!20] \p circle (1.8pt);
  \node[font=\scriptsize,violet!60!black] at (-2.5,2.0) {unseen spoofs};
  % arrows
  \draw[tflow] (2.35,0.2) -- (1.75,0.1);
  \node[font=\tiny,red!60!black,align=center] at (2.6,-0.15) {pushed out\\($m_1$)};
  \draw[flow,green!45!black] (0.62,0.62) -- (0.28,0.3);
  % center
  \fill[black] (0,0) circle (1.4pt);
  \node[font=\tiny] at (0.28,0.14) {$\mathbf{w}_0$};
\end{tikzpicture}
\caption{Geometry of the one-class contrastive objective. OC-Softmax compresses
bona-fide embeddings within an angular margin $m_0$ of the learnable direction
$\mathbf{w}_0$ and pushes spoofs beyond $m_1$. Because \emph{no} spoof-specific
structure is modeled inside the target region, unseen generators (violet) fall outside
it by default, unlike a binary decision boundary fit only to seen attacks.}
\label{fig:oneclass}
\end{figure}

\noindent\textbf{OC-Softmax.} The one-class softmax loss~\cite{zhang2021oneclass}
compresses bona-fide embeddings toward a learnable direction $\mathbf{w}_0$ while
pushing spoof embeddings away (Figure~\ref{fig:oneclass}). With
$\hat{\mathbf{z}} = \mathbf{z}/\lVert\mathbf{z}\rVert$,
$\hat{\mathbf{w}}_0 = \mathbf{w}_0/\lVert\mathbf{w}_0\rVert$, and
$\cos\theta = \hat{\mathbf{w}}_0^\top \hat{\mathbf{z}}$:
\begin{equation}
\mathcal{L}_{\text{OCS}} = \frac{1}{N}\sum_{i=1}^{N} \log\!\Big(1 +
  e^{\,s_0 (m_{y_i} - \cos\theta_i)(-1)^{y_i}}\Big),
\label{eq:ocs}
\end{equation}
where $s_0$ is a scale factor and $m_0 > m_1$ are asymmetric margins: bona-fide
embeddings must satisfy $\cos\theta > m_0$ (compactness) while spoof embeddings must
satisfy $\cos\theta < m_1$ (repulsion). At inference the score is simply
$s(\mathbf{x}) = -\cos\theta$, interpretable as angular distance from the bona-fide
prototype---a property we exploit for calibration (Section~\ref{ssec:calibration}).

\noindent\textbf{Supervised contrastive auxiliary.} A projection head
$g(\mathbf{z}) \in \R^{128}$ applies the SupCon loss~\cite{khosla2020supcon} within each
minibatch:
\begin{equation}
\mathcal{L}_{\text{SupCon}} = \sum_{i=1}^N \frac{-1}{|\mathcal{P}(i)|}
    \sum_{p \in \mathcal{P}(i)} \log \frac{
        \exp(g(\mathbf{z}_i)^\top g(\mathbf{z}_p) / \tau_c)
    }{ \sum_{j \neq i} \exp(g(\mathbf{z}_i)^\top g(\mathbf{z}_j) / \tau_c) },
\label{eq:supcon}
\end{equation}
where $\mathcal{P}(i)$ indexes same-label samples and $\tau_c$ is a temperature. Whereas
OC-Softmax constrains embeddings only relative to the single direction $\mathbf{w}_0$,
SupCon shapes the \emph{local neighborhood structure}, pulling bona-fide samples from
different speakers, channels, and augmentations together. The combination yields a
bona-fide manifold that is compact both globally and locally; ablation shows the two
terms are complementary (Section~\ref{ssec:ablation}). One-class variants with adaptive
centroids~\cite{kim2024oneclassknowledge} are orthogonal and could be integrated.

\subsection{Channel-Robustness Augmentation Curriculum}
\label{ssec:augmentation}

To address G2, each training waveform passes, with probability $p_{\text{aug}} = 0.7$,
through a randomly sampled distortion pipeline (Figure~\ref{fig:aug}):

\begin{figure}[t]
\centering
\resizebox{\columnwidth}{!}{%
\begin{tikzpicture}[node distance=3mm]
\node[io] (x) {$\mathbf{x}$};
\node[aug,right=4mm of x] (rb) {RawBoost\\LnL / ISD / SSI};
\node[aug,right=of rb] (codec) {codec chain\\MP3/AAC/Opus\\AMR/G.711};
\node[aug,right=of codec] (rir) {reverb\\(RIR conv.)};
\node[aug,right=of rir] (noise) {noise\\MUSAN\\0--20 dB};
\node[io,right=4mm of noise] (xa) {$\tilde{\mathbf{x}}$};
\draw[flow] (x) -- (rb);
\draw[flow] (rb) -- (codec);
\draw[flow] (codec) -- (rir);
\draw[flow] (rir) -- (noise);
\draw[flow] (noise) -- (xa);
\node[font=\tiny,gray,below=1.5mm of rb] {$p{=}0.5$};
\node[font=\tiny,gray,below=1.5mm of codec] {$p{=}0.5$};
\node[font=\tiny,gray,below=1.5mm of rir] {$p{=}0.3$};
\node[font=\tiny,gray,below=1.5mm of noise] {$p{=}0.4$};
\draw[flow,gray,dashed] (x.north) to[out=25,in=155] node[above,font=\tiny]{skip, $1-p_{\text{aug}}$} (xa.north);
\end{tikzpicture}%
}
\caption{Channel-robustness augmentation curriculum. Each stage is applied
independently with the indicated probability; the whole pipeline is skipped with
probability $1-p_{\text{aug}}$. Stage intensities are annealed from mild to full range
over the first 10 epochs.}
\label{fig:aug}
\end{figure}

\begin{enumerate}
    \item \textbf{RawBoost}~\cite{tak2022rawboost}: linear-nonlinear convolutive
        distortion, impulsive signal-dependent noise, and stationary signal-independent
        noise.
    \item \textbf{Codec chain:} one of MP3 (64--320~kbps), AAC (64--256~kbps), Opus
        (12--64~kbps), AMR-NB (4.75--12.2~kbps), or G.711 (A-law/$\mu$-law), applied via
        ffmpeg encode--decode round trips; with probability 0.15 two codecs are chained
        to simulate re-uploads.
    \item \textbf{Reverberation:} convolution with a measured room impulse response
        from the OpenSLR-28 collection~\cite{ko2017rir}.
    \item \textbf{Additive noise:} a MUSAN~\cite{snyder2015musan} noise, music, or
        babble clip at $\mathrm{SNR} \in [0, 20]$~dB.
\end{enumerate}

The pipeline is a \emph{curriculum}: stage intensities (bitrate floor, SNR floor,
RawBoost gain) are annealed linearly from mild to full range over the first 10 epochs,
which we found stabilizes the one-class objective---aggressive distortion from epoch~0
prevents the bona-fide region from forming. Augmentation is disabled at evaluation.

\subsection{Training and Inference Procedures}
\label{ssec:trainproc}

\begin{algorithm}[t]
\caption{AuralGuard training (one epoch)}
\label{alg:training}
\footnotesize
\KwIn{Training set $\mathcal{D}$; frozen SSL encoder $E_{\text{SSL}}$; trainable
  parameters $\Theta$ (layer attention $\boldsymbol{\alpha}$, View-B CNN, fusion,
  back-end, heads); epoch index $e$}
\ForEach{minibatch $\{(\mathbf{x}_i, y_i)\}_{i=1}^{B} \subset \mathcal{D}$}{
  $\tilde{\mathbf{x}}_i \leftarrow \textsc{AugmentCurriculum}(\mathbf{x}_i;\, e)$
    \tcp*{Sec.~\ref{ssec:augmentation}}
  $\mathbf{H}_i \leftarrow E_{\text{SSL}}(\tilde{\mathbf{x}}_i)$
    \tcp*{no grad, checkpointed}
  $\mathbf{F}_{A,i} \leftarrow \operatorname{Conv1d}\!\big(\textstyle\sum_\ell
    \pi_\ell \mathbf{h}_{\ell,i}\big)$ \tcp*{Eq.~\eqref{eq:layerattn}}
  $\mathbf{F}_{B,i} \leftarrow \operatorname{CNN}_{\text{art}}
    (\textsc{LFCC}\oplus\textsc{MGD}\oplus\textsc{CQTph}(\tilde{\mathbf{x}}_i))$\;
  $\mathbf{F}_i \leftarrow \textsc{GatedXAttn}(\mathbf{F}_{A,i}, \mathbf{F}_{B,i})$\;
  $\mathbf{z}_i \leftarrow \textsc{AASIST}(\mathbf{F}_i)$\;
  $\mathcal{L} \leftarrow \mathcal{L}_{\text{OCS}} + \lambda_c
    \mathcal{L}_{\text{SupCon}} + \lambda_r \mathcal{L}_{\text{ent}}$
    \tcp*{Eqs.~\eqref{eq:loss}--\eqref{eq:supcon}}
  $\Theta \leftarrow \textsc{AdamW}(\Theta, \nabla_\Theta \mathcal{L})$\;
}
\end{algorithm}

\noindent\textbf{Training.} Algorithm~\ref{alg:training} summarizes one epoch. Only
$\approx$4.6M of 321M total parameters receive gradients (Table~\ref{tab:budget});
activations of the frozen encoder are gradient-checkpointed so the model trains on a
single 24~GB GPU.

\noindent\textbf{Inference.} Clips longer than 4~s are scored with overlapping 4~s
windows (50\% hop); window scores are pooled by $\operatorname{logsumexp}$, which
approximates a soft maximum and reflects the forensic reality that a clip is synthetic
if \emph{any} segment is. The raw angular score is mapped to a probability via a
sigmoid with temperature $T$ fitted on the development set (temperature
scaling~\cite{guo2017calibration}); Section~\ref{ssec:calibration} shows that the
one-class score requires markedly less correction than BCE logits. The decision
threshold $\tau$ is selected on the development set for a target false-alarm rate.

\subsection{Parameter and Computation Budget}
\label{ssec:budget}

% [DUMMY] — parameter/MACs/RTF values to be replaced by measured numbers
\begin{table}[t]
\centering
\caption{Parameter and computation budget (4~s input). Trainable parameters exclude
  the frozen WavLM-Large encoder. RTF = real-time factor, lower is better.
  \emph{Values are placeholders to be replaced by measured numbers.}}
\label{tab:budget}
\footnotesize
\setlength{\tabcolsep}{4pt}
\begin{tabular}{lrr}
\toprule
Component & Params & GMACs \\
\midrule
WavLM-Large (frozen)            & 315.5M & 71.4 \\
Layer attention + projection    & 0.16M  & 0.5  \\
View-B feature CNN              & 1.1M   & 1.8  \\
Gated cross-attention fusion    & 0.53M  & 0.7  \\
AASIST back-end (GATv2)         & 2.6M   & 1.2  \\
OC-Softmax + SupCon heads       & 0.06M  & $<0.1$ \\
\midrule
Total (trainable / all)         & 4.6M / 320.1M & 75.7 \\
\midrule
RTF (GPU, RTX 3090)             & \multicolumn{2}{r}{0.06} \\
RTF (CPU, 8 threads)            & \multicolumn{2}{r}{0.41} \\
\bottomrule
\end{tabular}
\end{table}

Table~\ref{tab:budget} details the budget. The multi-view additions (View~B + fusion)
contribute only 1.6M parameters and $\approx$3\% of total MACs on top of the standard
SSL+AASIST recipe, so the generalization gains reported in Section~\ref{sec:results}
are not attributable to added capacity.
